% generated by 作成スクリプト/make_steps.py; do not edit
\begin{calc}[手順 1：ずらす向きを決める式]
比較対象からのずれを $\bm a = \bm w - \bm w_{\text{mkt}}$ とおく。3 銘柄とも $w_i>0$ で「空売りなし」は効いていないので、効いている制約は「合計 1」と「TE $=$ 上限」の 2 つ。この 2 つに乗数 $\lambda,\nu$ を付けたラグランジュ関数
\[ \mathcal L = \bm\mu^\top\bm w - \tfrac{\delta}{2}\bm w^\top\Sigma\bm w - \lambda(\mathbf 1^\top\bm w - 1) - \nu\bigl((\bm w-\bm w_{\text{mkt}})^\top\Sigma(\bm w-\bm w_{\text{mkt}}) - (L/\sqrt{252})^2\bigr) \]
を $\bm w$ で微分して 0 とおく：
\begin{align*}
\bm\mu - \delta\Sigma\bm w - \lambda\mathbf 1 - 2\nu\Sigma(\bm w-\bm w_{\text{mkt}}) &= \bm 0\\
\bm\mu - \delta\Sigma\bm w_{\text{mkt}} - \delta\Sigma\bm a - \lambda\mathbf 1 - 2\nu\Sigma\bm a &= \bm 0 &&(\bm w = \bm w_{\text{mkt}} + \bm a \text{ を代入})\\
(\delta+2\nu)\,\Sigma\bm a &= \bm\mu - \bm\pi - \lambda\mathbf 1 &&(\delta\Sigma\bm w_{\text{mkt}} = \bm\pi)\\
\bm a &= k\,\Sigma^{-1}(\bm\mu-\bm\pi) - k\lambda\,\Sigma^{-1}\mathbf 1 = k\,(\bm g - \lambda\bm e) &&(k = 1/(\delta+2\nu))
\end{align*}
ここで $\bm g = \Sigma^{-1}(\bm\mu-\bm\pi)$、$\bm e = \Sigma^{-1}\mathbf 1$。$\lambda$ は「ずれの合計 $a_1+a_2+a_3=0$」から、$k$ は「TE $=$ 上限」から決まる。
\end{calc}
\begin{calc}[手順 2：$\Sigma^{-1}$、$\bm g$、$\bm e$]
\[ \Sigma^{-1} = \begin{pmatrix}271.7 & -145.6 & -76.9\\-145.6 & 2{,}480.6 & -373.8\\-76.9 & -373.8 & 1{,}901.6\end{pmatrix},\qquad \bm\mu-\bm\pi = (0.0002895,\ 0.0000022,\ -0.0000857) \]
\begin{align*}
g_1 &= 271.7\times0.0002895 + (-145.6)\times0.0000022 + (-76.9)\times(-0.0000857)\\ &= 0.0787 + (-0.0003) + 0.0066 = 0.0849\\
g_2 &= (-145.6)\times0.0002895 + 2{,}480.6\times0.0000022 + (-373.8)\times(-0.0000857)\\ &= (-0.0422) + 0.0056 + 0.0320 = -0.0046\\
g_3 &= (-76.9)\times0.0002895 + (-373.8)\times0.0000022 + 1{,}901.6\times(-0.0000857)\\ &= (-0.0223) + (-0.0008) + (-0.1630) = -0.1861\\
e_1 &= 271.7 + (-145.6) + (-76.9) = 49.2\\
e_2 &= (-145.6) + 2{,}480.6 + (-373.8) = 1{,}961.1\\
e_3 &= (-76.9) + (-373.8) + 1{,}901.6 = 1{,}451.0
\end{align*}
\end{calc}
\begin{calc}[手順 3：$\lambda$（ずれの合計を 0 にする）]
\begin{align*}
\lambda &= \frac{g_1+g_2+g_3}{e_1+e_2+e_3} = \frac{0.0849 + (-0.0046) + (-0.1861)}{49.2 + 1{,}961.1 + 1{,}451.0} = \frac{-0.1058}{3{,}461.4} = -0.0000306\\
g_1 - \lambda e_1 &= 0.0849 - (-0.0000306)\times49.2 = 0.0849 - (-0.0015) = 0.0864\\
g_2 - \lambda e_2 &= -0.0046 - (-0.0000306)\times1{,}961.1 = -0.0046 - (-0.0599) = 0.0554\\
g_3 - \lambda e_3 &= -0.1861 - (-0.0000306)\times1{,}451.0 = -0.1861 - (-0.0443) = -0.1418
\end{align*}
向き $\bm g-\lambda\bm e = (0.0864, 0.0554, -0.1418)$ の合計は $0.0864 + 0.0554 + (-0.1418) = 0.0000$。
\end{calc}
\begin{calc}[手順 4：$k$（TE を上限にちょうど合わせる）]
向き $\bm v=\bm g-\lambda\bm e$ の「事前の TE の 2 乗」$\bm v^\top\Sigma\bm v$ を、まず $\Sigma\bm v$、次に $\bm v^\top(\Sigma\bm v)$ の順に計算する：
\begin{align*}
(\Sigma\bm v)_1 &= 0.0038780\times0.0864 + 0.0002590\times0.0554 + 0.0002077\times(-0.1418)\\ &= 320.084\times10^{-6}\\
(\Sigma\bm v)_2 &= 0.0002590\times0.0864 + 0.0004327\times0.0554 + 0.0000955\times(-0.1418)\\ &= 32.793\times10^{-6}\\
(\Sigma\bm v)_3 &= 0.0002077\times0.0864 + 0.0000955\times0.0554 + 0.0005530\times(-0.1418)\\ &= -55.176\times10^{-6}\\
\bm v^\top\Sigma\bm v &= (0.0864\times320.084 + 0.0554\times32.793 + (-0.1418)\times(-55.176))\times10^{-6}\\ &= 0.00003730\\
k &= \frac{L/\sqrt{252}}{\sqrt{0.00003730}} = \frac{0.001588}{0.006108} = 0.26002
\end{align*}
上限がなければ $k = 1/\delta = 0.40$ で、そのときの TE は $0.006108\times0.40\times\sqrt{252} = 3.88\%$。これが上限 $2.52\%$ を超えるので、上限が効いて $k$ が $0.40$ から $0.26002$ に縮む（$\nu = (1/k-\delta)/2 = (3.8459 - 2.5)/2 = 0.6729 > 0$）。
\end{calc}
\begin{calc}[手順 5：比率]
\begin{align*}
\textsf{SANA}:\ a_1 &= 0.26002\times0.0864 = 0.0225, & w_1 &= 0.3000 + 0.0225 = \mathbf{0.3225}\\
\textsf{AIV}:\ a_2 &= 0.26002\times0.0554 = 0.0144, & w_2 &= 0.2500 + 0.0144 = \mathbf{0.2644}\\
\textsf{JACK}:\ a_3 &= 0.26002\times(-0.1418) = -0.0369, & w_3 &= 0.4500 + (-0.0369) = \mathbf{0.4131}
\end{align*}
最適化のソフト（Clarabel）の答え：$0.3225, 0.2644, 0.4131$。手で解いた値と一致する。合計 $0.3225 + 0.2644 + 0.4131 = 1.0000$。
\end{calc}
\begin{calc}[検算：事前の TE]
\[ \bm a^\top\Sigma\bm a = \sum_{i,j} a_i\Sigma_{ij}a_j = 2.5221\times10^{-6},\quad \sqrt{2.5221\times10^{-6}}\times\sqrt{252} = 0.001588\times15.8745 = 2.52\% \]
上限 $2.52\%$ と一致する。
\end{calc}
